\section{Every minimally compensated vertical-lattice derivative}
\label{sec:lattice}

The vertical-lattice report \cite{RepoVertical} proves a normally
convergent sampling formula after Euler--Maclaurin compensation, and
computes a single extra Bernoulli term at the first negative-even
derivative. Its Questions 1 and 6 ask for all higher minimally
compensated derivatives and for the periodic filters that remove their
contacts. We answer both questions. The answer is a single polynomial
in the logarithm of the cutoff, specified by an elementary generating
function. In particular, higher derivative order does not require
additional moments of a periodic filter.

\subsection{The inherited sampling identity and its exact domain}

Let \(h>0\), \(\tau=2\pi/h\), and \(\Ree a>0\). All powers use the
analytic logarithm on the right half-plane. Write
\[
 w_n=a+\ii hn,\qquad q=e^{-\tau},\qquad z=e^{-\tau a}.
\]
The generalized Stieltjes constants have the convention
\begin{equation}\label{eq:lattice-stieltjes}
 \zeta(1+u,w)=\frac1u+
       \sum_{m=0}^{\infty}\frac{(-1)^m}{m!}\gamma_m(w)u^m,
 \qquad \gamma_0(w)=-\psi(w).
\end{equation}
We use \(B_1=-1/2\). For \(M\geq0\), define
\begin{align}
 R_M(s,w)&=\zeta(s,w)-\frac{w^{1-s}}{s-1}-\frac12w^{-s}
 -\sum_{j=1}^{M}\frac{B_{2j}}{(2j)!}(s)_{2j-1}w^{1-s-2j},
 \label{eq:root-remainder}\\
 {\cal L}_{\nu}(q,z)&=\sum_{k=1}^{\infty}
             \frac{z^k}{k^\nu(1-q^k)},\label{eq:root-lambert}\\
 {\cal B}_M(s;a,\tau)&={\cal L}_{1-s}(q,z)
       -\frac12\Li_{1-s}(z)-\frac1\tau\Li_{2-s}(z)\notag\\
 &\hspace{7mm}
       -\sum_{j=1}^{M}\frac{B_{2j}\tau^{2j-1}}{(2j)!}
                    \Li_{2-s-2j}(z),\label{eq:root-bracket}\\
 {\cal F}_M(s;a,h)&=\frac{\tau^s}{\Gamma(s)}{\cal B}_M(s;a,\tau).
 \label{eq:root-spectral}
\end{align}
The removable value at \(s=1\) is understood in \(R_M\).
Both \(R_M(s,w)\) and \({\cal F}_M(s;a,h)\) are entire in \(s\).

For completeness, the previously established sampling formula is
\begin{equation}\label{eq:root-master}
 \sum_{n\in\mathbb Z}R_M(s,a+\ii hn)={\cal F}_M(s;a,h),
 \qquad \Ree s>-2M.
\end{equation}
It includes all fixed order derivatives on compact subsets of that
strict half-plane. One direct proof applies Poisson periodization to
\[
 \frac{t^{s-1}e^{-at}}{\Gamma(s)}
 \left(\frac1{1-e^{-t}}-\frac1t-\frac12
       -\sum_{j=1}^{M}\frac{B_{2j}}{(2j)!}t^{2j-1}\right)
 \quad(t>0),
\]
extended by zero for \(t\leq0\). The density tends to zero at the
origin on the stated sufficient half-plane, and its
Fourier coefficients are absolutely summable there. The standard
Hurwitz expansion \cite[Section 25.11]{DLMF} gives the required
normal convergence. We credit \eqref{eq:root-master} to the preceding
report; it is the starting point, not a new claim.

\subsection{A universal power-lattice cutoff lemma}

The following elementary lemma isolates the whole obstruction.

\begin{lemma}\label{lem:power-cutoff}
For \(a,h\) as above and \(u\) in a sufficiently small complex
neighborhood of zero, put
\[
 P_K(u)=\sum_{n=-K}^{K}(a+\ii hn)^{-1-u},\qquad
 P_\infty^{\mathrm{an}}(u)=
       \frac{\tau^{1+u}}{\Gamma(1+u)}\Li_{-u}(e^{-\tau a}).
\]
Then, locally uniformly in \(u\) and \(a\),
\begin{equation}\label{eq:power-cutoff}
 P_K(u)=P_\infty^{\mathrm{an}}(u)
    +\frac{2}{h}\frac{\sin(\pi u/2)}{u}
                   e^{-u\log(hK)}
    +O\!\left(K^{-1-\Ree u}\right).
\end{equation}
The sine quotient is filled by its removable value. Derivatives of
order \(j\) of the error at zero are
\(O(K^{-1}(1+\log K)^j)\).
\end{lemma}

\begin{proof}
For \(\Ree u>0\), ordinary periodization of
\(t^u e^{-at}/\Gamma(1+u)\) gives
\(\sum_n(a+\ii hn)^{-1-u}=P_\infty^{\mathrm{an}}(u)\).
The omitted paired tail has the expansion
\[
 (a+\ii hn)^{-1-u}+(a-\ii hn)^{-1-u}
 =-2\sin(\pi u/2)(hn)^{-1-u}+O(n^{-2-\Ree u}).
\]
Since
\(\sum_{n>K}n^{-1-u}=K^{-u}/u+O(K^{-1-\Ree u})\),
subtracting the tail proves \eqref{eq:power-cutoff} in that half-plane.
To see the uniform continuation through \(u=0\), use the exact tail
expression
\begin{align*}
 &( \ii h)^{-1-u}\zeta(1+u,K+1-\ii a/h)\\
 &\hspace{12mm}
 +(-\ii h)^{-1-u}\zeta(1+u,K+1+\ii a/h).
\end{align*}
Its two residues cancel at \(u=0\). Euler--Maclaurin applied to
these two Hurwitz functions proves the stated expansion on a full
neighborhood. Differentiating that expansion introduces only powers
of \(\log K\). The sector and parameter estimates are uniform on
compact subsets of \(\Ree a>0\).
\end{proof}

At \(u=0\), the lemma gives the familiar distinction
\[
 \lim_{K\to\infty}P_K(0)
 =\frac{\tau}{2}\coth\frac{\tau a}{2}
 =\tau\Li_0(e^{-\tau a})+\frac{\tau}{2}.
\]
The extra \(\tau/2\) is the zeroth member of the entire cutoff
generating function in \eqref{eq:power-cutoff}.

\subsection{The all-order contact polynomial}

For an even nonnegative integer \(N\), define
\begin{equation}\label{eq:root-beta}
 M=N/2,\qquad
 \beta_N=\frac{B_{N+2}}{(N+1)(N+2)},\qquad
 A_N(u)=\prod_{j=1}^{N}\left(1-\frac uj\right).
\end{equation}
An empty product equals one. Define the polynomial
\begin{equation}\label{eq:contact-polynomial}
 \boxed{\displaystyle
 C_{N,d}(X)=
  \frac{\tau\beta_N}{2}\,d!\,[u^d]\!
  \left\{uA_N(u)\frac{\sin(\pi u/2)}{\pi u/2}e^{-Xu}\right\},
  \qquad d\geq1.}
\end{equation}
Its degree is \(d-1\), with nonzero leading coefficient
\((\tau\beta_N/2)d(-1)^{d-1}\).

\begin{theorem}[Complete minimal-contact law]\label{thm:minimal-contact}
For every even \(N\geq0\), every \(d\geq1\), and \(\Ree a>0\),
\begin{equation}\label{eq:minimal-contact}
 \boxed{\displaystyle
 \sum_{n=-K}^{K}
   \left.\partial_s^dR_{N/2}(s,a+\ii hn)\right|_{s=-N}
 =
 {\cal F}_{N/2}^{(d)}(-N;a,h)
 +C_{N,d}\bigl(\log(hK)\bigr)
 +O\!\left(\frac{(1+\log K)^{d-1}}K\right).}
\end{equation}
Here and below the superscript on \({\cal F}\) denotes the order
derivative. The error is uniform for \(a\) in compact subsets of
the right half-plane.

For odd \(N\geq1\), put \(M=(N+1)/2\). Then all the corresponding
derivative sums are absolutely convergent, and
\begin{equation}\label{eq:odd-all-jets}
 \sum_{n\in\mathbb Z}
 \left.\partial_s^dR_M(s,a+\ii hn)\right|_{s=-N}
 ={\cal F}_M^{(d)}(-N;a,h).
\end{equation}
\end{theorem}

\begin{proof}
For even \(N\), the difference of two successive compensations is
exactly
\begin{equation}\label{eq:one-more-subtraction}
 R_{N/2}(-N+u,w)-R_{N/2+1}(-N+u,w)
 =\beta_N uA_N(u)w^{-1-u}.
\end{equation}
Indeed,
\[
 \frac{B_{N+2}}{(N+2)!}(-N+u)_{N+1}
       =\beta_N uA_N(u).
\]
The oversubtracted series is normally convergent in
\(\Ree u>-2\). Its omitted tail after \(d\) derivatives at zero is
\(O(K^{-2}(1+\log K)^{d-1})\), since its leading coefficient
contains the factor \(u\). In the initial half-plane \(\Ree u>0\),
\eqref{eq:root-master} and \eqref{eq:one-more-subtraction} also imply
\[
 {\cal F}_{N/2}(-N+u)-{\cal F}_{N/2+1}(-N+u)
       =\beta_NuA_N(u)P_\infty^{\mathrm{an}}(u).
\]
Both sides continue holomorphically. Apply Lemma~\ref{lem:power-cutoff}
to the truncated version of \eqref{eq:one-more-subtraction}. Its
explicit cutoff contribution is
\[
 \frac{2\beta_N}{h}A_N(u)\sin(\pi u/2)e^{-u\log(hK)}
 =\frac{\tau\beta_N}{2}\,
       uA_N(u)\frac{\sin(\pi u/2)}{\pi u/2}e^{-u\log(hK)}.
\]
Taking the \(d\)th coefficient proves the theorem and the stated
error. For odd \(N\), the minimally chosen \(M\) already satisfies
\(-N>-2M\), so \eqref{eq:odd-all-jets} is an immediate application
of normal convergence in \eqref{eq:root-master}.
\end{proof}

This proof also explains why merely differentiating the first-order
Bernoulli correction is insufficient: the factor
\(A_N(u)\sin(\pi u/2)/(\pi u/2)\) itself contributes higher
coefficients.

\begin{corollary}[Two explicit cutoff conventions]\label{cor:cutoff-coordinate}
The constant term of the cutoff expansion in powers of \(\log K\)
is
\begin{equation}\label{eq:index-finite-part}
 \FP_{\log K}\sum_{n\in\mathbb Z}
     \partial_s^dR_{N/2}(-N,a+\ii hn)
 ={\cal F}_{N/2}^{(d)}(-N;a,h)+C_{N,d}(\log h).
\end{equation}
The constant term when the large variable is instead the physical
radius \(R=hK\) is
\begin{equation}\label{eq:radius-finite-part}
 \FP_{\log R}\sum_{|hn|\leq R}
     \partial_s^dR_{N/2}(-N,a+\ii hn)
 ={\cal F}_{N/2}^{(d)}(-N;a,h)+C_{N,d}(0).
\end{equation}
Only positive powers of the specified logarithm are discarded in
either definition. Their difference is exactly
\(C_{N,d}(\log h)-C_{N,d}(0)\).
\end{corollary}

\begin{proof}
Theorem~\ref{thm:minimal-contact} supplies a polynomial plus a term
that tends to zero. Taking its constant coefficient in the indicated
variable proves both formulas. Replacing \(R/h\) by its integer part
changes the asymptotic by a quantity tending to zero.
\end{proof}

\subsection{Harmonic and zeta coefficients}

Let \(Y_j\) denote the complete exponential Bell polynomials,
defined by
\[
 \exp\left(\sum_{r=1}^{\infty}x_r\frac{u^r}{r!}\right)
   =\sum_{j=0}^{\infty}Y_j(x_1,\ldots,x_j)\frac{u^j}{j!}.
\]
The Euler product for sine and the finite product in
\eqref{eq:root-beta} give
\begin{align}
 x_1&=-X-H_N,\notag\\
 x_r&=-(r-1)!\left(H_N^{(r)}
                +\mathbf1_{2\mid r}\,2^{1-r}\zeta(r)\right),
                  \qquad r\geq2,\label{eq:contact-bell}\\
 C_{N,d}(X)&=\frac{\tau\beta_N}{2}\,d\,
       Y_{d-1}(x_1,\ldots,x_{d-1}).\label{eq:contact-bell-value}
\end{align}
Thus the contact coefficients involve finite harmonic numbers and
even zeta values only; no unspecified Laurent coefficient remains.
Writing \(L=X+H_N\) and
\(V=H_N^{(2)}+\zeta(2)/2\), the first four are
\begin{align}
 C_{N,1}&=\frac{\tau\beta_N}{2},&
 C_{N,2}&=-\tau\beta_N L,\notag\\
 C_{N,3}&=\frac{3\tau\beta_N}{2}(L^2-V),&
 C_{N,4}&=2\tau\beta_N
       \bigl(-L^3+3LV-2H_N^{(3)}\bigr).
 \label{eq:first-four-contacts}
\end{align}
For \(N=0\), in particular,
\begin{equation}\label{eq:gamma-cubic-contact}
 C_{0,1}=\frac{\tau}{24},\qquad
 C_{0,2}(X)=-\frac{\tau X}{12},\qquad
 C_{0,3}(X)=\frac{\tau}{8}
                    \left(X^2-\frac{\pi^2}{12}\right).
\end{equation}
The constant \(-\tau\pi^2/96\) at the third derivative is missed
if the sine factor is dropped.

For explicit evaluation of the analytic term in
\eqref{eq:minimal-contact}, define \(E_j^{(N)}\) by
\begin{equation}\label{eq:root-E}
 \exp\left((\log\tau+\gamma-H_N)u+
   \sum_{r=2}^{\infty}
       \frac{(-1)^{r+1}\zeta(r)-H_N^{(r)}}r u^r\right)
   =\sum_{j=0}^{\infty}E_j^{(N)}\frac{u^j}{j!}.
\end{equation}
The reciprocal-Gamma factorization then gives, for every \(M\),
\begin{equation}\label{eq:root-analytic-jets}
 {\cal F}_M^{(d)}(-N;a,h)
 =(-1)^NN!\tau^{-N}d
   \sum_{j=0}^{d-1}\binom{d-1}{j}E_j^{(N)}
       \partial_s^{d-1-j}{\cal B}_M(-N;a,\tau).
\end{equation}
This coefficient formula is inherited from \cite{RepoVertical};
the new result supplies its exact minimal-cutoff interpretation.
All order derivatives on the right are normally convergent
polylogarithm or Lambert--polylogarithm derivatives.

\subsection{Periodic filters: the mean is the only obstruction}

Let \(c_n\) have period \(p\) and mean
\(\bar c=p^{-1}\sum_{r=0}^{p-1}c_r\). Define
\begin{equation}\label{eq:filtered-analytic}
 {\cal F}_{M,c}(s;a,h)=
    \sum_{r=0}^{p-1}c_r\,
       {\cal F}_M(s;a+\ii hr,ph).
\end{equation}

\begin{theorem}[All-order filter criterion]\label{thm:filter-all}
For even \(N\) and \(d\geq1\),
\begin{equation}\label{eq:weighted-contact}
 \sum_{n=-K}^{K}c_n\partial_s^dR_{N/2}(-N,a+\ii hn)
 =
 {\cal F}_{N/2,c}^{(d)}(-N;a,h)
 +\bar c\,C_{N,d}(\log(hK))+o(1).
\end{equation}
If \(\bar c=0\), the positive and negative halves both converge
ordinarily, locally uniformly in the order parameter near \(-N\);
every derivative equals the corresponding derivative of
\eqref{eq:filtered-analytic}. No finite part is needed.
The condition \(\bar c=0\) is necessary and sufficient for all
contacts to vanish. There are no additional discrete-moment
conditions at higher derivative order.
\end{theorem}

\begin{proof}
Decompose \(c_n=\bar c+(c_n-\bar c)\). The constant part is
Theorem~\ref{thm:minimal-contact}. Partial sums of the mean-zero
part are bounded. Summation by parts, using
\[
 \partial_s^dR_{N/2}(-N,a+\ii hn)
     =O\!\left(\frac{(1+\log|n|)^{d-1}}{|n|}\right)
\]
and its first difference
\(O((1+\log|n|)^{d-1}/n^2)\), proves convergence of each half.
The same argument before differentiation is locally uniform in
\(\Ree(s+N)>-1\). In \(\Ree(s+N)>0\) the sums equal
\eqref{eq:filtered-analytic} by absolute convergence and residue-class
splitting; analytic continuation proves equality near \(s=-N\).
Necessity follows already from
\(C_{N,1}=\tau\beta_N/2\neq0\).
\end{proof}

For the alternating filter this yields the particularly simple
all-order identity
\begin{equation}\label{eq:alternating-all-jets}
 \sum_{n\in\mathbb Z}(-1)^n
      \partial_s^dR_{N/2}(-N,a+\ii hn)
 =
 2{\cal F}_{N/2}^{(d)}(-N;a,2h)
   -{\cal F}_{N/2}^{(d)}(-N;a,h).
\end{equation}
The two half-sums specify the summation convention; absolute
convergence is not asserted. For every two shifts \(a,b\) in the
right half-plane, the relative identity
\begin{equation}\label{eq:relative-lattice}
 \begin{split}
 &\sum_{n\in\mathbb Z}c_n
 \left\{\partial_s^dR_{N/2}(-N,a+\ii hn)
       -\partial_s^dR_{N/2}(-N,b+\ii hn)\right\}\\
 &\hspace{10mm}
 ={\cal F}_{N/2,c}^{(d)}(-N;a,h)
       -{\cal F}_{N/2,c}^{(d)}(-N;b,h)
 \end{split}
\end{equation}
is absolutely convergent for every bounded periodic \(c\).
The leading reciprocal-logarithmic term is independent of the
horizontal shift, so subtraction improves it to
\(O(n^{-2}(1+\log|n|)^{d-1})\).

\subsection{Gamma and all Stieltjes antiderivatives}

At \(N=0\), write \(D_d(w)=\partial_s^dR_0(0,w)\). Direct
differentiation of the fixed subtractions gives
\begin{equation}\label{eq:Dd-explicit}
 D_d(w)=\zeta^{(d)}(0,w)
       +d!\,w\sum_{j=0}^{d}\frac{(-\Log w)^j}{j!}
       -\frac12(-\Log w)^d.
\end{equation}
The first instance is the usual Stirling remainder,
\[
 D_1(w)=\log\Gamma(w)-(w-\tfrac12)\Log w+w-\tfrac12\log(2\pi).
\]
Its vertical Gamma product belongs to the published work of
Arabi Ardehali and Rosengren \cite{ArdehaliRosengren}, as the
preceding report correctly records. Our contribution here is the
explicit all-\(d\) minimal-cutoff and filter law.

For example,
\begin{equation}\label{eq:D2-cutoff}
 \begin{split}
 &\sum_{n=-K}^{K}
 \left\{\zeta''(0,w_n)+w_n\bigl(\Log^2w_n-2\Log w_n+2\bigr)
                       -\tfrac12\Log^2w_n\right\}\\
 &\hspace{9mm}
 ={\cal F}_0''(0;a,h)-\frac{\tau}{12}\log(hK)
            +O(K^{-1}\log K).
 \end{split}
\end{equation}
This identity is simultaneously a formula for compensated
antiderivatives of the first generalized Stieltjes constant:
expansion of
\(\partial_w\zeta(s,w)=-s\zeta(s+1,w)\) gives
\begin{equation}\label{eq:stieltjes-primitive-root}
 \partial_w\zeta^{(d)}(0,w)=(-1)^d d\,\gamma_{d-1}(w),
 \qquad d\geq1.
\end{equation}
Thus
\[
 {\mathfrak A}_m(w)=\frac{(-1)^{m+1}}{m+1}
       \left(\zeta^{(m+1)}(0,w)-\zeta^{(m+1)}(0,1)\right)
\]
is the uniquely anchored primitive
\(\mathfrak A_m'(w)=\gamma_m(w)\),
\(\mathfrak A_m(1)=0\). Theorem~\ref{thm:minimal-contact},
\eqref{eq:Dd-explicit}, and \eqref{eq:root-analytic-jets}
evaluate every corresponding compensated lattice sum using finite
harmonic/zeta coefficients and Lambert--polylogarithm order
derivatives. Formula \eqref{eq:relative-lattice} supplies an
ordinary absolutely convergent version whenever two anchors are
compared.

